\section{Exact depth bounds under the six standard argument changes}
\label{gorbit:sec:main}

The existing continuation proves that, at weight six, products and the
six transformations permuting $0,1,\infty$ leave exactly one formal
triple direction outside the span of doubles. It also observes that the
same weight-six span is obtained from all invertible linear changes of
the two differential forms. We extend both calculations to every weight
and depth. The equality of these two spans at weight six does not persist:
at weight seven their dimensions are respectively $12$ and $14$.

The height-one application belongs to the established theory of Nielsen
polylogarithms. Charlton--Gangl--Radchenko give its polynomial model and
depth bound in \cite[\S3.3, Theorem~7]{gorbit:CGR}; their Proposition~28
is exactly the weight-seven congruence below. We supply a formal
membership proof, an explicit weight-eight Lie certificate, and a fully
replayable product completion of that known congruence. The principal
extension here is the rank formula for the full binary shuffle quotient,
including every irreducible component and any finite set of directions.

All lower bounds in this section concern the specified formal shuffle
presentation. They do not assert numerical independence at $z=i$, and
they impose no obstruction on a proof using additional cyclotomic
letters, distribution, stuffle, or other functional equations.

\subsection{The precise presentation}

Let $V=\mathbb Qx\oplus\mathbb Qy$ and let $\operatorname{Sh}(V)$
be the shuffle algebra on binary words. Write $J$ for its augmentation
ideal, and set
\[
 Q_w=(J/J^2)_w,\qquad
 D_{w,d}=\operatorname{span}\{[v]: |v|=w,
                 \text{$v$ contains at most $d$ copies of $y$}\}.
\]
Here $[v]$ denotes the class modulo shuffle products of nonempty words.
We use the full binary shuffle algebra so that every linear letter
substitution acts immediately. At weights $w\ge2$, adding words ending
in $x$ does not enlarge the indecomposables: shuffle normalization with
$H_x(z)=\log z$ removes trailing $x$ letters.

The analytic interpretation is outer letter first:
\[
 x=\frac{dt}{t},\qquad y=\frac{dt}{1-t},\qquad
 H_{x^{s_1-1}y\cdots x^{s_r-1}y}(z)
   =\operatorname{Li}_{s_1,\ldots,s_r}(z,1,\ldots,1).
\]
Pure $x$ words have the standard logarithmic regularization. Products
are assigned the sum of the weights of their factors.

The two pullback matrices
\begin{equation}\label{gorbit:eq:matrices}
 R=\begin{pmatrix}0&-1\\-1&0\end{pmatrix},\qquad
 M=\begin{pmatrix}1&0\\1&-1\end{pmatrix}
\end{equation}
act on columns. Thus $R(x)=-y$, $R(y)=-x$, while
$M(x)=x+y$, $M(y)=-y$. They correspond to $z\mapsto1-z$ and
$z\mapsto z/(z-1)$. The group
\[
 \mathfrak S=\{I,R,M,RM,MR,RMR\}
\]
has six elements. It sends the line $\mathbb Qx$ to precisely the
three lines $\mathbb Qx$, $\mathbb Qy$, and $\mathbb Q(x+y)$.
Define
\[
 E^{(6)}_{w,d}=\sum_{g\in\mathfrak S}gD_{w,d},\qquad
 E^{(\mathrm{lin})}_{w,d}
    =\sum_{g\in\operatorname{GL}_2(\mathbb Q)}gD_{w,d}.
\]
The second space is a formal enlargement. A general matrix in
$\operatorname{GL}_2$ need not arise from an argument change preserving
the three-puncture alphabet.

Complementation changes the initial endpoint from zero to one.
Its analytic realization therefore includes products with regularized
endpoint constants. Whenever an analytic congruence is written below,
the quotient is by the weight-$w$ part of the product ideal generated by
two positive-weight hyperlogarithms or endpoint constants, and by pure
endpoint constants of weight $w$. No product is discarded in a stated
full identity.

\subsection{A formula depending only on the number of projective directions}

For $w\ge2$ and $0\le j\le w$, put
\begin{equation}\label{gorbit:eq:witt}
 \ell_{w,j}=\frac1w\sum_{k\mid\gcd(w,j)}
       \mu(k)\binom{w/k}{j/k},\qquad
 m_{w,j}=\ell_{w,j}-\ell_{w,j-1}\quad
       (1\le j\le\lfloor w/2\rfloor).
\end{equation}
In particular, $\ell_{w,0}=0$. The $\ell_{w,j}$ are the binary
multigraded Witt numbers already used in the manuscript. The integers
$m_{w,j}$ are nonnegative.

\begin{theorem}[Finite-direction depth-orbit formula]
\label{gorbit:thm:finite}
Let $w\ge2$ and $1\le d\le\lfloor w/2\rfloor$. Let
$\mathcal G$ be any nonempty finite collection of invertible rational
letter maps, and suppose that $\{[gx]:g\in\mathcal G\}$ contains
exactly $k$ distinct points of $\mathbb P(V)$. Then
\begin{equation}\label{gorbit:eq:finite}
 \boxed{\quad
 \dim\left(\sum_{g\in\mathcal G}gD_{w,d}\right)
 =\sum_{j=1}^{d}m_{w,j}
       \min\{w-2j+1,\ k(d-j+1)\}.
 \quad}
\end{equation}
Consequently
\begin{align}
 \dim E^{(6)}_{w,d}
   &=\sum_{j=1}^{d}m_{w,j}
          \min\{w-2j+1,3(d-j+1)\},\label{gorbit:eq:six}\\
 \dim E^{(\mathrm{lin})}_{w,d}
   &=\sum_{j=1}^{d}m_{w,j}(w-2j+1).
          \label{gorbit:eq:linear}
\end{align}
Thus the finite-direction dimension depends only on $k$, not on the
positions of the distinct rational directions.
\end{theorem}

\begin{proof}
We separate the representation-theoretic and interpolation arguments.
The shuffle polynomial theorem and the standard Lyndon Lie basis imply
that the weight-$j$ subspace of $Q_w$, where weight here counts $y$
letters, has dimension $\ell_{w,j}$. Equivalently, its two-variable
character is the usual Witt character
\[
 \chi_{Q_w}(a,b)=\frac1w\sum_{k\mid w}\mu(k)(a^k+b^k)^{w/k}.
\]
These are classical free-Lie and shuffle facts; see
\cite{gorbit:Reutenauer}. They may also be obtained by
counting primitive binary necklaces and choosing their Lyndon
representatives, as in the existing manuscript.

Extend scalars to $\mathbb C$ temporarily. Polynomial representations
of $\operatorname{GL}_2$ are completely reducible
\cite{gorbit:Etingof}, and their homogeneous
irreducibles of degree $w$ have the form
\[
 W_{w,j}=\operatorname{Sym}^{w-2j}V\otimes(\det V)^j.
\]
This representation has one-dimensional weight spaces with $y$ counts
$j,j+1,\ldots,w-j$. Comparing the character coefficients successively
for $0\le j\le w/2$ gives
\begin{equation}\label{gorbit:eq:decomposition}
 Q_w\simeq\bigoplus_{j=1}^{\lfloor w/2\rfloor}
        W_{w,j}^{\oplus m_{w,j}}.
\end{equation}
Indeed $\ell_{w,j}=\sum_{a\le j}m_{w,a}$, so subtraction gives
exactly \eqref{gorbit:eq:witt}. This also proves the nonnegativity.
Ranks of rational matrices are unchanged by the scalar extension.

Fix one copy of $W_{w,j}$, with $j\le d$, and write
$n=w-2j$ and $r=d-j$. Its intersection with $D_{w,d}$ is
\[
 x^{n-r}\operatorname{Sym}^{r}V\otimes(\det V)^j.
\]
This is the span of $x^n,x^{n-1}y,\ldots,x^{n-r}y^r$, with the
determinant factor understood. A map $g$ sends it to
\[
 (gx)^{n-r}\operatorname{Sym}^{r}V\otimes(\det V)^j.
\]
The determinant contributes a nonzero scalar and does not affect its
span. Components with $j>d$ have no intersection with $D_{w,d}$.

It remains to calculate the dimension of the sum of the spaces
$l_a^{n-r}\operatorname{Sym}^rV$ for $k$ distinct lines $[l_a]$.
Associate to a linear functional $\lambda$ on $\operatorname{Sym}^nV$
the homogeneous polynomial
\[
 P_\lambda(s,t)=\lambda((sx+ty)^n).
\]
This identifies the dual with binary forms of degree $n$, since every
binomial coefficient involved is nonzero. The functional annihilates
$l^{n-r}\operatorname{Sym}^rV$ precisely when $P_\lambda$ has a zero
of multiplicity at least $r+1$ at the projective point $[l]$:
differentiate $(l+\varepsilon m)^n$ through order $r$ for any
independent $m$. For $k$ distinct points, these conditions mean
divisibility by the product of their $k$ distinct linear vanishing
forms, each to the power $r+1$. The annihilator therefore has dimension
\[
 \max\{0,n-k(r+1)+1\}.
\]
Subtracting from $n+1$ gives $\min\{n+1,k(r+1)\}$.
This is a direct homogeneous interpolation proof; no generic-position
assumption beyond distinctness of the points is needed.

Summing independently over the copies in \eqref{gorbit:eq:decomposition}
proves \eqref{gorbit:eq:finite}. The six matrices have three directions,
which proves \eqref{gorbit:eq:six}. Finally, sufficiently many distinct
rational directions fill each of the finitely many symmetric-power
components, proving \eqref{gorbit:eq:linear} for
$\operatorname{GL}_2(\mathbb Q)$ itself.
\end{proof}

\subsection{Consequences for the depth-two search}

Since $m_{w,1}=1$ and
$m_{w,2}=\lfloor(w-1)/2\rfloor-1$, the double case simplifies to
\begin{align}
 \dim E^{(6)}_{w,2}
 &=\min\{w-1,6\}
  +\left(\left\lfloor\frac{w-1}{2}\right\rfloor-1\right)
                           \min\{w-3,3\},\label{gorbit:eq:double-six}\\
 \dim E^{(\mathrm{lin})}_{w,2}
 &=(w-3)\left\lfloor\frac{w-1}{2}\right\rfloor+2,
 \qquad w\ge4.\label{gorbit:eq:double-linear}
\end{align}
In particular, the six-argument rank is
$3\lfloor(w-1)/2\rfloor+3$ for $w\ge7$.

\begin{center}
\begin{tabular}{c|rrrr}
$w$&$\dim Q_w$&$\dim E^{(6)}_{w,2}$&
 $\dim E^{(\mathrm{lin})}_{w,2}$&$\dim(Q_w/E^{(6)}_{w,2})$\\\hline
4&3&3&3&0\\
5&6&6&6&0\\
6&9&8&8&1\\
7&18&12&14&6\\
8&30&12&17&18\\
9&56&15&26&41\\
10&99&15&30&84\\
11&186&18&42&168\\
12&335&18&47&317
\end{tabular}
\end{center}

This table makes two different limitations visible. Components with
$j>d$ cannot be reached by any linear-letter orbit of the depth-$d$
subspace. Within components with $j\le d$, a finite set of argument
directions can still fail to span the whole symmetric power. The
weight-six calculation sees only the first limitation. Both are already
present at weight seven.

For fixed $d$ and $w\ge3d+1$, every minimum in
\eqref{gorbit:eq:six} selects its second entry, and telescoping gives
\begin{equation}\label{gorbit:eq:stable}
 \dim E^{(6)}_{w,d}=3\sum_{j=1}^{d}\ell_{w,j}.
\end{equation}
Hence, as $w\to\infty$ with $d$ fixed,
\[
 \dim E^{(6)}_{w,d}\sim\frac{3w^{d-1}}{d!},\qquad
 \dim E^{(\mathrm{lin})}_{w,d}\sim\frac{w^d}{d!},\qquad
 \dim Q_w\sim\frac{2^w}{w}.
\]
The estimates follow directly from the $k=1$ term in
\eqref{gorbit:eq:witt}; all other divisor terms have smaller polynomial
degree for fixed $j$, and the usual binary Witt sum gives the last
estimate. Thus a fixed depth budget and a fixed set of standard arguments
span a polynomially growing subspace of an exponentially growing formal
space.

\subsection{A sharp height-one threshold}

The general dimensions hide a particularly simple invariant family.
Write
\[
 T_{w,h}=[x^{w-h}y^h],\qquad 1\le h\le w-1.
\]
Its analytic representative is
$\operatorname{Li}_{w-h+1,\{1\}^{h-1}}(z)$, or $S_{w-h,h}(z)$ in
Nielsen notation. The uniform bound
$\lceil(w-1)/3\rceil=\lfloor(w+1)/3\rfloor$ is
\cite[Theorem~7]{gorbit:CGR}. The following individual-word formulation
records exact membership in our specified formal presentation, using
the same polynomial structure; it is not a claim of a new analytic
Nielsen depth bound.

\begin{theorem}[Exact formal depth of height-one words]
\label{gorbit:thm:height}
The smallest integer $d\ge1$ for which
$T_{w,h}\in E^{(6)}_{w,d}$ is
\begin{equation}\label{gorbit:eq:height-depth}
 \boxed{\qquad
 d_{\min}(w,h)=\min\left\{h,w-h,
                       \left\lceil\frac{w-1}{3}\right\rceil\right\}.
 \qquad}
\end{equation}
In particular,
$[x^{w-d-1}y^{d+1}]\notin E^{(6)}_{w,d}$ whenever $w\ge3d+2$.
The threshold is sharp: at $w=3d+1$ every height-one word lies in
$E^{(6)}_{w,d}$.
\end{theorem}

\begin{proof}
Let $L=L(V)$ be the free Lie algebra inside the tensor algebra, and put
$L'=[L,L]$, $L''=[L',L']$. Coefficient pairing identifies degree-$w$
Lie polynomials with the dual of $Q_w$: Lie polynomials annihilate
shuffle products, and the standard Lyndon Lie basis has the required
dimension.
For a letter matrix $g$ and a Lie polynomial $P$, the coefficient
pairing satisfies
\[
 \langle P,gq\rangle=\langle g^{\mathsf T}P,q\rangle,
\]
where $g^{\mathsf T}$ denotes the induced substitution by the transpose
letter matrix. This is an adjoint identity; the contragredient group
action on the dual would instead use $g^{-\mathsf T}$. In particular,
stability of a Lie subspace under all letter substitutions makes its
annihilator in $Q_w$ stable as well.

Every nonzero monomial in an element of $L'$ contains both letters.
Consequently no concatenation of two such monomials is an ordered word
$x^a y^b$. Thus $T_{w,h}$ annihilates $L''$. The degree-$w$ part of
$L'/L''$ is naturally
\begin{equation}\label{gorbit:eq:metabelian}
 (L'/L'')_w\simeq\operatorname{Sym}^{w-2}V\otimes\det V.
\end{equation}
For completeness, $\operatorname{ad}_x$ and $\operatorname{ad}_y$
commute on $L'/L''$, since their commutator is
$\operatorname{ad}_{[x,y]}$, which vanishes there. Jacobi's identity
therefore spans the quotient by
\[
 \operatorname{ad}_x^a\operatorname{ad}_y^b[x,y],\qquad a+b=w-2.
\]
These $w-1$ elements are independent. Map $x$ and $y$ to upper
triangular matrices over $\mathbb Q[s,t]$ by
\[
 x\longmapsto\begin{pmatrix}s&0\\0&0\end{pmatrix},\qquad
 y\longmapsto\begin{pmatrix}t&1\\0&0\end{pmatrix}.
\]
The displayed brackets have upper-right entries $s^{a+1}t^b$, which
are distinct monomials. Their derived ideal is abelian, so this map
factors through the metabelian quotient. Naturality of $[x,y]$ under
letter changes proves \eqref{gorbit:eq:metabelian}.

It follows that every $T_{w,h}$ lies in the unique $j=1$ summand of
\eqref{gorbit:eq:decomposition}. Put $n=w-2$. In that summand the
depth-$d$ orbit under the six matrices is
\[
 x^{n-d+1}\operatorname{Sym}^{d-1}V
 +y^{n-d+1}\operatorname{Sym}^{d-1}V
 +(x+y)^{n-d+1}\operatorname{Sym}^{d-1}V,
\]
with its determinant factor suppressed. By the interpolation proof,
its annihilator is zero if $n<3d$, and otherwise is
\begin{equation}\label{gorbit:eq:ann-height}
 \Delta(s,t)^d\operatorname{Sym}^{n-3d}(\mathbb Q^2)^*,
 \qquad \Delta(s,t)=st(s-t).
\end{equation}

The vector $T_{w,h}$ is a nonzero vector in its one-dimensional
$y$-weight-$h$ space. The corresponding coefficient in the dual binary
form is its $t^{h-1}$ coefficient, up to a nonzero scalar. To see the
normalization explicitly, set $A=\operatorname{ad}_x^{w-1}y$ and let
$\delta$ be the derivation $\delta(x)=y$, $\delta(y)=0$.
For $0\le r\le n$, the Lie polynomial
$\delta^r A/(n)_r$ corresponds to $s^{n-r}t^r$; here
$(n)_r=n(n-1)\cdots(n-r+1)$ and $(n)_0=1$. In this Lie realization,
annihilating $D_{w,d}$ means divisibility by $t^d$. Applying the adjoints
of $R$ and $M$ requires also $s^d$ and $(s-t)^d$, respectively. This
recovers \eqref{gorbit:eq:ann-height} and fixes its sign convention.
The coefficient of $\delta^r A/(n)_r$ at
$x^{w-h}y^h$ is zero unless $r=h-1$, when it equals $(-1)^r$.
Indeed
\[
 r!\sum_{k=0}^{r}(-1)^k\binom{w-1}{k}
   =(-1)^r r!\binom{w-2}{r}=(-1)^r(n)_r.
\]

If $n\ge3d$, the possible $t$ exponents with nonzero coefficient in
\eqref{gorbit:eq:ann-height} are precisely
$d,d+1,\ldots,n-d$. Every such exponent occurs: multiply
$\Delta^d$ by a suitable monomial of degree $n-3d$, and use the
nonzero binomial coefficients in $(s-t)^d$.
Thus the annihilator vanishes on $T_{w,h}$ exactly when
\[
 h\le d,\qquad\text{or}\qquad h\ge w-d,
 \qquad\text{or}\qquad w\le3d+1.
\]
Solving these three alternatives for the smallest $d$ gives
\eqref{gorbit:eq:height-depth}.
\end{proof}

The formula combines the original and complementary depth bounds with
the established uniform Nielsen bound. Our proof also tests exact
nonmembership against all binary words of the permitted depth, rather
than just against the height-one subfamily. At weight four it gives
depth one for every height-one word; at weights six and seven it gives
depth two for the middle height-one words.

\subsection{Explicit primitive obstructions in every weight}

The preceding proof gives actual rational certificates of nonmembership.
For $w\ge3d+2$, set $n=w-2$, $A=\operatorname{ad}_x^{w-1}y$, and
\begin{equation}\label{gorbit:eq:omega-general}
 \Omega_{w,d}=
 \sum_{r=0}^{d}(-1)^r\binom dr
          \frac{\delta^{d+r}A}{(n)_{d+r}}.
\end{equation}
It corresponds to the dual form
$s^{n-3d}\Delta(s,t)^d$, so it annihilates
$E^{(6)}_{w,d}$. Being a Lie polynomial, it also annihilates every
shuffle product of two nonempty words. Its coefficient at
$x^{w-d-1}y^{d+1}$ is $(-1)^d$, proving that the target class survives.

The weight-eight double certificate is especially short:
\begin{equation}\label{gorbit:eq:omega8}
 \boxed{\quad
 \Omega_{8,2}
 =\frac{\delta^4-6\delta^3+12\delta^2}{360}
       \operatorname{ad}_x^7y,
 \qquad
 \langle\Omega_{8,2},x^5y^3\rangle=1.
 \quad}
\end{equation}
Hence $\operatorname{Li}_{6,1,1}(z)$ cannot be reduced in the formal
six-argument double presentation. This is an explicit nonzero
functional. The corresponding Nielsen obstruction for
$S_{5,3}$ is already discussed in \cite[\S9, preconsiderations]{gorbit:CGR};
our contribution here is the normalized Lie certificate with its exact
full-binary verification, not priority for the nonreduction phenomenon.

\subsection{An exact product completion of the known weight-seven identity}

Write $F_{a,b}(z)=\operatorname{Li}_{a,b}(z,1)$ and
$F_{a,b,c}(z)=\operatorname{Li}_{a,b,c}(z,1,1)$. Put
$u=1-z$ and $v=z/(z-1)$.

\begin{proposition}[Charlton--Gangl--Radchenko, Proposition~28]
\label{gorbit:prop:weight7}
Modulo positive-weight products and endpoint constants,
\begin{align}
 F_{5,1,1}(z)\equiv{}&-3\operatorname{Li}_7(z)+2F_{6,1}(z)
       +2\operatorname{Li}_7(u)-F_{6,1}(u)\notag\\
       &-3\operatorname{Li}_7(v)+F_{6,1}(v).
 \label{gorbit:eq:weight7-congruence}
\end{align}
This is \cite[Proposition~28]{gorbit:CGR} in our decreasing-index
notation. A complete product lift is given by
\eqref{gorbit:eq:weight7-full} below.
\end{proposition}

To state the full identity without hiding a long product polynomial,
let $a=x^6y$, $b=x^5y^2$, $t=x^4y^3$, and define the explicit
weight-seven word polynomial
\begin{equation}\label{gorbit:eq:residual7}
 \mathcal R=t+3a-2b-2R(a)+R(b)+3M(a)-M(b).
\end{equation}
Expand it as $\mathcal R=\sum_w c_w w$. For this particular word
polynomial define
\begin{equation}\label{gorbit:eq:product7}
 \mathcal P_{\mathcal R}(z)
 =\sum_w c_w\sum_{k=2}^{7}\frac{(-1)^k}{k}
       \sum_{\substack{w=w_1\cdots w_k\\|w_i|>0}}
                      \prod_{i=1}^{k}H_{w_i}(z).
\end{equation}
All factors have strictly smaller weight than seven. This is a finite,
explicit expression; the accompanying JSON lists its $864$ nonzero
commutatively collected product monomials. Its exact shuffle
re-expansion is $\mathcal R$.

Put $\alpha=-\log(1-z)$ and
\begin{align}
 C_7(z)&=\zeta(7)+\sum_{k=1}^{5}
                   \frac{(-1)^k\alpha^k}{k!}\zeta(7-k),\notag\\
 C_{6,1}(z)&=\zeta(6,1)+\sum_{k=1}^{4}
                   \frac{(-1)^k\alpha^k}{k!}\zeta(6-k,1).
 \label{gorbit:eq:endpoint7}
\end{align}
On the principal domain
\[
 \mathcal U=\mathbb C\setminus
                  \bigl(( -\infty,0]\cup[1,\infty)\bigr),
\]
the full identity is
\begin{align}
 F_{5,1,1}(z)={}&-3\operatorname{Li}_7(z)+2F_{6,1}(z)
       +2\operatorname{Li}_7(u)-F_{6,1}(u)\notag\\
 &-3\operatorname{Li}_7(v)+F_{6,1}(v)
       +\mathcal P_{\mathcal R}(z)-2C_7(z)+C_{6,1}(z).
 \label{gorbit:eq:weight7-full}
\end{align}
In particular, it may be specialized at $z=i$. The factors with trailing
$x$ in \eqref{gorbit:eq:product7} have their ordinary shuffle
normalization in $\log z$, so no new undefined functions occur.

\begin{proof}
There are eighteen weight-seven binary Lyndon words. Pairing
\eqref{gorbit:eq:residual7} with their eighteen standard Lie brackets
gives zero in every coordinate. This exact integer calculation is
reconstructed by the supplied verifier. Since these brackets form the
dual of $Q_7$, it proves $\mathcal R\in J^2$.

We explain why \eqref{gorbit:eq:product7} is a complete product
correction, rather than just a certificate in a quotient. In a connected
commutative graded Hopf algebra, the convolution logarithm of the
identity is
\[
 e=\log^\star(\mathrm{id})
   =\sum_{k\ge1}\frac{(-1)^{k-1}}k
                       (\mathrm{id}-\eta\varepsilon)^{\star k}.
\]
It annihilates products of positive-degree elements. One elementary
proof uses the convolution powers $\mathrm{id}^{\star m}$ for
nonnegative integers $m$: they are algebra homomorphisms because the
target is commutative. In every bounded degree these powers are
polynomial in $m$. Their homomorphism identity therefore holds for a
formal parameter; taking its coefficient of degree one in that
parameter gives the infinitesimal-character identity for $e$.

The coproduct of the shuffle Hopf algebra is deconcatenation. Thus
$e(\mathcal R)=0$ is exactly the formal identity saying that the
shuffle re-expansion of \eqref{gorbit:eq:product7} equals
$\mathcal R$. Evaluation gives
$H_{\mathcal R}(z)=\mathcal P_{\mathcal R}(z)$. Independently of
this general argument, the verifier expands all $864$ products and
checks that their exact rational residual is zero.

For $0<z<1$, Landen pullback gives
$H_{M(a)}(z)=\operatorname{Li}_7(v)$ and
$H_{M(b)}(z)=F_{6,1}(v)$. Path composition at the regularized endpoint
one gives
\[
 H_{R(a)}(z)=\operatorname{Li}_7(1-z)-C_7(z),\qquad
 H_{R(b)}(z)=F_{6,1}(1-z)-C_{6,1}(z).
\]
For example, the suffixes of $a=x^6y$ contribute $\zeta(7-k)$;
the complementary prefix $x^k$ becomes $(-y)^k$ and contributes
$(-1)^k\alpha^k/k!$. The suffix $y$ has regularized endpoint value
zero. The same calculation for $b=x^5y^2$ gives $\zeta(6-k,1)$;
the suffixes $y$ and $y^2$ again have zero regularized constant.
These observations give exactly \eqref{gorbit:eq:endpoint7}.
Substituting the four identities into $H_{\mathcal R}=\mathcal P_{
\mathcal R}$ proves \eqref{gorbit:eq:weight7-full} on $(0,1)$.
Analytic continuation proves it on $\mathcal U$. Discarding precisely
the named product and endpoint terms gives
\eqref{gorbit:eq:weight7-congruence}.
\end{proof}

\subsection{Independent exact verification and remaining questions}

The program \texttt{verify\_depth\_orbits.py} was written independently
of the previous weight-six checker. It constructs binary Lyndon Lie
polynomials by recursive commutators, applies the displayed integer
letter substitutions, and computes matrix ranks over $\mathbb Q$ using
integer elimination with gcd normalization. It imports no numerical
polylogarithm evaluator and no previous relation matrix.

At weights $2$ through $10$, it checks all $25$ pairs
$(w,d)$ with $1\le d\le\lfloor w/2\rfloor$, both for the six maps
and for enough distinct rational linear directions to fill the full
linear orbit. It additionally checks $34$ depth-two cases involving
one through five irregularly spaced rational directions. Every rank
agrees with \eqref{gorbit:eq:finite}. A further $155$ exact membership
checks verify the height-one threshold for every relevant height and
depth through weight ten. At weight eight the explicit
obstruction is checked on all $222$ transformed binary words of depth
at most two, including trailing-zero words, and on all $1792$
nonempty binary shuffle products. Its target pairing is exactly one.
The weight-seven product certificate is independently re-expanded
term by term with exact fractions.

These finite checks audit the conventions and examples. The
all-weight statements follow from the representation and interpolation
proofs, without extrapolation from a table. The Witt character,
complete reducibility, binary-form interpolation, and convolution
logarithm are classical ingredients. The finite-direction formula treats
the full binary quotient; the height-one component recovers established
Nielsen structure. The independently constructed exact certificates and
product lift make these statements directly auditable. None of the
height-one reductions is presented as a new identity, and no claim of
worldwide priority is made for the general rank formula.

Several next questions are now sharply posed. First, one can construct
compact product corrections for other height-one reductions predicted
by \eqref{gorbit:eq:height-depth}, replacing the deliberately uniform
all-cuts expression by smaller identities. Second, the irreducible
components with $j>d$ isolate the exact missing representations for
future argument alphabets. Third, it remains to determine which of the
formal obstructions survive after specialization to Gaussian or
Eisenstein points and after adding their full colored relations.
Finally, the conjectured short $S_6$ identity remains a separate
mixed-color problem: none of the binary rank or height-one theorems
proves or refutes it.

% Bibliography entries are supplied in depth_orbits_bibliography.tex.
% All CGR theorem/proposition numbers above refer to arXiv:1908.04770v1.
% The inspected repository source is
% Analysis/Polylogarithms/docs/reports/finite-golden-uniform-continuation/article/depth.tex
% at commit 3447bc59b78d138a7f0d536b66ac6e5cc5d5e49f.
